\newcommand{\NVIDIA}{%
  \tikz[baseline=(char.base)]{
    \node[
      circle,
      fill=white!60!white,   % 深绿色底色（可调整）
      text=green,            % 白色字母
      inner sep=1.2pt,
      draw=black,            % 黑色边框（可选）
      font=\sffamily\bfseries\fontsize{6}{7}\selectfont
    ] (char) {N};
  }%
}

\newcommand{\Google}{%
  \tikz[baseline=(char.base)]{
    \node[
      circle,
      fill=white!60!white,   % 深绿色底色（可调整）
      text=blue,            % 白色字母
      inner sep=0.6pt,
      draw=black,            % 黑色边框（可选）
      font=\sffamily\bfseries\fontsize{8}{7}\selectfont
    ] (char) {G};
  }%
}

\newcommand{\HUAWEI}{%
  \tikz[baseline=(char.base)]{
    \node[
      circle,
      fill=white!60!white,   % 深绿色底色（可调整）
      text=red,            % 白色字母
      inner sep=0.6pt,
      draw=black,            % 黑色边框（可选）
      font=\sffamily\bfseries\fontsize{8}{7}\selectfont
    ] (char) {H};
  }%
}

\newcommand{\AMD}{%
  \tikz[baseline=(char.base)]{
    \node[
      circle,
      fill=white!60!white,   % 深绿色底色（可调整）
      text=black,            % 白色字母
      inner sep=0.6pt,
      draw=black,            % 黑色边框（可选）
      font=\sffamily\bfseries\fontsize{8}{7}\selectfont
    ] (char) {A};
  }%
}

\newcommand{\speedup}{%
  \tikz[baseline=(char.base)]{
    \node[
      circle,
      fill=black!60!black,   % 深绿色底色（可调整）
      text=white,            % 白色字母
      inner sep=0.6pt,
      draw=blue,            % 黑色边框（可选）
      font=\sffamily\bfseries\fontsize{8}{7}\selectfont
    ] (char) {S};
  }%
}

\newcommand{\correctness}{%
  \tikz[baseline=(char.base)]{
    \node[
      circle,
      fill=black!60!black,   % 深绿色底色（可调整）
      text=white,            % 白色字母
      inner sep=0.6pt,
      draw=blue,            % 黑色边框（可选）
      font=\sffamily\bfseries\fontsize{8}{7}\selectfont
    ] (char) {C};
  }%
}


\newcommand{\efficiency}{%
  \tikz[baseline=(char.base)]{
    \node[
      circle,
      fill=black!90!black,   % 深绿色底色（可调整）
      text=white,            % 白色字母
      inner sep=0.6pt,
      draw=blue,            % 黑色边框（可选）
      font=\sffamily\bfseries\fontsize{8}{7}\selectfont
    ] (char) {E};
  }%
}

\newcommand{\fastp}{%
  \tikz[baseline=(char.base)]{
    \node[
      circle,
      fill=red!100!red,   % 深绿色底色（可调整）
      text=white,            % 白色字母
      inner sep=0.6pt,
      draw=black,            % 黑色边框（可选）
      font=\sffamily\bfseries\fontsize{8}{7}\selectfont
    ] (char) {f};
  }
}

\newcommand{\perf}{%
  \tikz[baseline=(char.base)]{
    \node[
      circle,
      fill=red!100!red,   % 深绿色底色（可调整）
      text=white,            % 白色字母
      inner sep=0.6pt,
      draw=black,            % 黑色边框（可选）
      font=\sffamily\bfseries\fontsize{8}{7}\selectfont
    ] (char) {P};
  }%
}

\newcommand{\similarity}{%
  \tikz[baseline=(char.base)]{
    \node[
      circle,
      fill=red!100!red,   % 深绿色底色（可调整）
      text=white,            % 白色字母
      inner sep=0.6pt,
      draw=black,            % 黑色边框（可选）
      font=\sffamily\bfseries\fontsize{8}{7}\selectfont
    ] (char) {S};
  }%
}

This chapter focuses on the systematic benchmarking of kernel generation, and provides a structured overview of representative evaluation benchmarks, including both evaluation metrics and benchmark datasets, to lay a solid foundation for subsequent method comparison and performance analysis.


\begin{table*}[ht]
\centering
% \begin{tabularx}{\textwidth}{XcccX}
\begin{threeparttable}
\begin{tabular}{p{3.2cm}p{1cm}p{1.2cm}p{1.3cm}p{8.8cm}}
\hline
\textbf{Name} & \textbf{Time} & \textbf{Metrics} & \textbf{Hardware} & \textbf{Description}   \\ \hline
ParEval             &  01/2024  & \correctness \  \speedup \  \efficiency   & \NVIDIA \AMD   &  420 expert-selected tasks across 12 algorithmic domains for benchmarking general parallel code generation.    \\
KernelBench & 02/2025  & \correctness \  \fastp     & \NVIDIA  & 250 PyTorch-to-CUDA kernel generation tasks, curated from popular GitHub repositories and official PyTorch operators, for evaluating AI/DL kernel generation.    \\
TritonBench & 02/2025  & \correctness \speedup \ \similarity  \  \efficiency\tnote{*} & \NVIDIA    & TritonBench evaluates Triton kernel generation via two subsets: 184 high-level kernels sourced from popular GitHub projects (TritonBench-G) and 166 fusion tasks derived from diverse PyTorch operators with different frequencies of usage (TritonBench-T).   \\
MultiKernel-Bench & 07/2025 & \correctness \  \speedup   & \HUAWEI \ \NVIDIA \  \Google   & 285-task benchmark across 14 operator categories for multi-platform DL kernel synthesis.  \\
TritonBench-revised \& ROCm Benchmark & 07/2025  & \correctness \ \speedup   & \AMD   & An AMD GPU-centric evaluation dataset comprising 30 expert-verified ROCm kernels and an adapted version of TritonBench-G, specifically optimized for AMD GPU performance benchmarking.     \\
Robust-kbench  & 09/2025  & \correctness \ \speedup  & \NVIDIA   & A robustness-focused benchmark featuring 9 specialized deep learning task categories, derived by refining and extending KernelBench.  \\
BackendBench & 09/2025 & \correctness \ \speedup & \NVIDIA & A rigorous evaluation framework that enforces PyTorch's official core library standards on 271 operators. Current use cases primarily leverage NVIDIA CUDA and Triton, yet the architecture remains backend-agnostic. \\
CUDAEval      & 10/2025  & \correctness \  \speedup   & \NVIDIA   & Leveraging 313 curated tasks from the the Stack v2 to benchmark the efficacy of reasoning transfer in CUDA code optimization.   \\
FlashInfer-Bench & 01/2026  & \correctness \  \fastp     & \NVIDIA  & Provide a unified schema describing kernel definitions, workloads, implementations, and evaluations, including eight representative kernel types used in LLM inference. \\
\hline
% \end{tabularx}
\end{tabular}
\begin{tablenotes}
\item[*] Efficiency here is defined as the ratio of the operator's measured throughput to the theoretical maximum performance.
\end{tablenotes}
\end{threeparttable}
\caption{\texorpdfstring{Benchmark datasets for kernel generation and optimization. Metrics: \protect\correctness{} Correctness, \protect\speedup{} Speedup, \protect\efficiency{} Efficiency, \protect\fastp{} $fast_p$, \protect\perf{} Perf, \protect\similarity{} Similarity. Hardware Platforms: \protect\NVIDIA{} NVIDIA GPUs, \protect\HUAWEI{} HUAWEI NPUs, 
\protect\Google{} Google TPUs, \protect\AMD{} AMD GPUs.}{Benchmark datasets for code generation and optimization in kernel generation. 
Metrics: Correctness, Speedup, Efficiency, $fast_p$, Perf, Similarity. Hardware Platforms: NVIDIA GPUs, HUAWEI NPUs, Google TPUs, AMD GPUs.}}
\label{tab:benchmarks}
\end{table*}

% TritonGym         & 10/2025  & \correctness \  \perf      & \NVIDIA   & Focus on derivative and synthetic tasks by mutating operator semantics to create Out-of-Distribution (OOD) kernels and introducing domain-specific language extension tasks that are rarely encountered in standard datasets.  \\

\subsection{Metrics}
Several factors should be considered when evaluating the performance of the operator implementation: correctness, efficiency, and so on. 
To build a comprehensive evaluation, existing benchmarks generally adopt execution-based unit tests, where the generated kernels will be compared with the standard implementations of CUDA/PyTorch. Given the instability of operator generation, each testing task usually involves multiple evaluations across $k$ random samples among $n$ times of generation.

\textbf{Correctness} primarily includes two aspects based on difficulty: (1) successful compilation and (2) consistency with the reference in multiple input-output comparisons.
Among various metrics used in code generation, \textit{pass@}$k$ is widely chosen, which calculates the probability that at least one correct implementation is generated in $k$ trials. The standard estimator is defined as:
\begin{equation}
\text{pass@}k \triangleq \mathbb{E}\left[1-\binom{n-c}{k}
/\binom{n}{k}\right],
\end{equation}
where the expectation is taken over kernel tasks and prompts, $c$ is the number of correct kernel implementations. 

\textbf{Efficiency} is another principal goal that kernel evaluation focuses on.
Speedup@$k$ measures how much faster a generated implementation is compared with the baselines by calculating
\begin{equation}
\text{speedup@}k \triangleq \mathbb{E}\left[\sum_{j=1}^n \left(\binom{j-1}{k-1}T^{\mathrm{base}}\right)/\left(\binom{n}{k}T_j\right)\right],
\end{equation}
where $T_j$ is the running time of the $j$-th generated implementation while $T^{\mathrm{base}}$ is the time consumed by the baseline. Note that the implementations are sorted by their performance, i.e., $T_1$ corresponds to the slowest and $T_n$ to the fastest.

% % robustness包括兼容各个平台、各种语言、测试边界条件
% \textbf{Compatibility} should also be considered when evaluating operator generation techniques. 
% Specifically, since operator generation relies on different hardware platforms and involves multiple low-source programming languages, LLMs tend to exhibit significant performance variation.
% For example, LLM is a good at writing GPU operator but fails in generating correct operators when facing to NPU tasks.
% To thoroughly assess robustness, many benchmarks strive to create more diverse testing tasks and testing environments, detailed in the following subsections.

In addition, Efficiency@$k$ refers to how effectively the generated operators utilize computation resources during execution, and Compatibility is considered when evaluating operator generation techniques across different hardware platforms or languages. Combined metrics are also used to evaluate multiple aspects of performance. For example, \textit{Perf$@$K} measures how close the best result from $K$ generated kernels is to a human expert performance. The \textit{$fast_p$} jointly evaluates the functional correctness and runtime performance of generated kernels. \textit{Similarity} uses 4 items (n-gram, weighted n-gram, syntax and dataflow) to measure the similarity between the generated code and the reference code.

\subsection{Benchmark Datesets}

As summarized in Table~\ref{tab:benchmarks}, kernel benchmarks are evolving from simple, single-platform evaluations toward comprehensive, real-world and generalized operator evaluation. We observe the following three key trends.

\paragraph{Metrics.} Moving beyond basic correctness and raw speedup (ParEval~\cite{pareval2024}, KernelBench), recent suites adopt composite objectives. Examples include efficiency metrics in TritonBench~\cite{Li_2025_TritonBench} and robustness assessments in Robust-kbench.

\paragraph{Hardware.} Evaluation is expanding beyond NVIDIA exclusivity. Compared to early benchmarks such as ParEval and KernelBench exclusively targeting NVIDIA GPUs, MultiKernelBench~\cite{wen2025multikernelbench} integrates HUAWEI NPUs and Google TPUs, while TritonBench-revised targets AMD GPUs. Additionally, NPUEval~\cite{kalade2025npueval} specifically targets power-sensitive kernels of neural processing units.

\paragraph{Content.} Workloads are shifting from generic algorithms to production-grade traces. KernelBench and TritonBench emphasize real-world PyTorch-to-CUDA or Triton kernel generation curated from popular GitHub repositories and The Stack v2. FlashInfer-Bench~\cite{xing2026flashinfer} standardizes 1,600 real-world LLM serving workloads, and BackendBench~\cite{saroufim2025backendbench} targets complex edge cases. 

% A reliable benchmark depends not only on well-designed evaluation metrics, but also, and often more critically, on the representativeness, diversity, and controllability of its test data. As summarized in Table~\ref{tab:benchmarks}, we systematically review the major benchmark suites used for evaluating automated kernel generation. Collectively, these benchmarks exhibit consistent trends, including a shift from basic to composite metrics, from single-GPU to multi-platform hardware support, and from limited kernels to diverse, real-world and generalized operator evaluation. \\
% $\bullet$ \textbf{Key metrics} Regarding evaluation metrics, benchmark design evolves from basic correctness-oriented evaluation toward composite metrics that capture multiple optimization objectives. Early benchmarks primarily rely on functional correctness and raw speedup, as exemplified by ParEval~\cite{pareval2024} and KernelBench. In contrast, more recent benchmarks adopt richer metric sets, including efficiency and similarity in TritonBench~\cite{Li_2025_TritonBench}, \textit{Perf}@$K$ metrics in TritonGym~\cite{TritonGym_2025}, and robustness-oriented speedup evaluation in Robust-kbench. This progression reflects the inherently multi-objective nature of kernel generation and optimization.\\
% $\bullet$ \textbf{Hardware platforms} From the hardware platform perspective, benchmarks evolve from being tightly coupled to a single GPU vendor toward increasingly comprehensive multi-platform coverage. As shown in Table~\ref{tab:benchmarks}, early benchmarks such as ParEval, KernelBench, and TritonBench exclusively target NVIDIA GPUs. Later efforts, exemplified by MultiKernelBench~\cite{wen2025multikernelbench}, explicitly incorporate heterogeneous accelerators spanning NVIDIA GPUs, HUAWEI NPUs, and Google TPUs, while ROCm triton benchmark and TritonBench-revised~\cite{Wang_2025_GEAK} further extend evaluation to AMD GPUs. And NPUEval~\cite{kalade2025npueval} proposes a benchmark for writing and evaluating neural processing units kernels in power-sensitive devices.
% \\
% $\bullet$ \textbf{Evaluation content} In terms of evaluation content, benchmarks progressively move from general parallel computation and limited kernels toward diverse, real-world and generalized generation tasks. Specifically, ParEval focuses on 420 expert-selected parallel tasks across 12 algorithmic domains, while KernelBench and TritonBench emphasize real-world PyTorch-to-CUDA or Triton kernel generation curated from popular GitHub repositories and The Stack v2. FlashInfer-Bench~\cite{xing2026flashinfer} provides a standardized framework by introducing a unified schema for eight representative kernel types that cover the core components of modern LLM inference, and curates a rigorous dataset of 1,600 workloads collected directly from SGLang serving traces, capturing actual usage patterns across both short and long sequences. BackendBench~\cite{saroufim2025backendbench} utilizes comprehensive test suites like OpInfo to ensure production-grade robustness across complex edge cases and data types. TritonGym further expands the scope by introducing mutated operator semantics, out-of-distribution kernels, and DSL extension tasks that are rarely encountered in standard datasets to evaluate the generalization ability. 